\SourceLeaf{114}
\pdfbookmark[1]{Göttingische gelehrte Anzeigen. 23. Dezember 1815.}{anzeige1815}
\begin{center}
\rule{0.65\linewidth}{0.4pt}\par
Göttingische gelehrte Anzeigen. 23. Dezember 1815.\par
\rule{0.65\linewidth}{0.4pt}
\end{center}

Am 7. Dezember legte Herr Professor \textsc{Gauss} der Königlichen Gesellschaft eine Vorlesung mit folgendem Titel vor:

\begin{center}
\emph{Ein weiterer neuer Beweis des Satzes, dass jedes Polynom in einer Variablen in reelle Faktoren ersten oder zweiten Grades zerlegt werden kann.}
\end{center}

Der hier ausgesprochene Satz, der wichtigste in der Theorie der Gleichungen, hat bekanntlich mehrere der bedeutendsten Mathematiker auf vielfältige Weise beschäftigt. Die wichtigsten Versuche, einen vollkommen strengen Beweis dafür zu geben, von \textsc{d'Alembert}, \textsc{Euler}, \textsc{Foncenex} und \textsc{Lagrange}, hat der Verfasser der vorliegenden Abhandlung bereits in einer vor sechzehn Jahren erschienenen Schrift zusammengestellt und geprüft. So sehr der Scharfsinn anzuerkennen ist, der jene Arbeiten auszeichnet, lässt sich doch nicht leugnen, dass sie alle mehr oder weniger Lücken lassen, die ihre Beweiskraft zerstören. Obwohl \textsc{Lagrange}s tiefgehende Untersuchungen die meisten jener Mängel behoben haben, lassen sich weder die Bemühungen dieses großen Mathematikers noch die scharfsinnige Art, wie \textsc{Laplace} den Gegenstand später behandelt hat, von dem Hauptvorwurf freisprechen, der alle jene versuchten Beweise trifft. Er besteht darin, dass die Sache so behandelt wurde, als sei lediglich die \emph{Form} der Wurzeln zu bestimmen, während ihre \emph{Existenz} vorausgesetzt wurde, ohne sie zu beweisen,
\SourceLeaf{115}
eine Schlussweise, die gerade hier nur trügerisch und tatsächlich eine echte \emph{petitio principii} ist. Alle erwähnten Beweise sind rein analytisch, auch der von \textsc{d'Alembert}, obwohl er in einer geometrischen Einkleidung erscheint, die für die Sache selbst völlig gleichgültig ist. Man könnte sich daher fast zu dem Schluss verleiten lassen, jene Voraussetzung, deren Unzulässigkeit übrigens dem genannten Analysten selbst entgangen ist, habe sich bei einer analytischen Behandlung gar nicht vermeiden lassen.

Der Verfasser der vorliegenden Abhandlung hatte den Gegenstand in der oben erwähnten Schrift auf eine ganz andere Weise behandelt und einen äußerst einfachen neuen Beweis gegeben. Dessen Besonderheit ist, dass er sich zum Teil auf die Geometrie der Lage stützt und im Übrigen hinsichtlich der Strenge nichts zu wünschen übrigzulassen scheint. Zugleich hatte er aber schon damals erklärt, dass er es keineswegs für unmöglich halte, auf rein analytischem Weg zu einem vollkommen strengen Beweis zu gelangen, und sich die ausführliche Entwicklung seiner Ideen für eine andere Zeit vorbehalten. Andere Beschäftigungen hatten ihn bisher von diesem Gegenstand abgehalten: Die vorliegende Abhandlung soll jene Zusage erfüllen.

Was diesen neuen Beweis selbst betrifft, so steht er dem erwähnten früheren an Einfachheit und Kürze freilich weit nach; doch liegt dies in der Beschaffenheit der subtilen Untersuchung selbst. Es bleibt stets erfreulich, neben einem äußerst einfachen, aber zum Teil auf andersartigem Stoff beruhenden Beweis des wichtigen Satzes noch einen zweiten zu besitzen, der zwar länger und kunstvoller, aber für sich ganz selbstständig ist. Eine Darstellung der Hauptideen des neuen Beweises würde, wenn sie auch nur einigermaßen befriedigend ausfallen sollte, weit mehr Raum beanspruchen, als der Zweck dieser Blätter erlaubt. Wir müssen daher diejenigen, die einen Hauptreiz der Mathematik in der vollkommenen Klarheit ihrer Wahrheiten und der höchsten Strenge ihrer Beweise finden, auf den bald erscheinenden dritten Band der Commentationes der Gesellschaft verweisen.
\begin{center}\rule{0.28\linewidth}{0.4pt}\end{center}
